\section{Hermetic Meditation}

\label{sec:HermeticMeditation}

Since Gornahoor has been so insistent on the necessity for spiritual practice, we will here provide an example of Hermetic meditation. This is based on suggestions from the writings of Valentin Tomberg, the foremost Hermetist of the 20th century. As he points out:

\begin{quotex}
There are no theories; there is only experience, including here the intellectual experience of arcana (mysteries) and symbols.

\end{quotex}
The process involves relaxing the mind, and then entering into the arcana using the active imagination, almost as an eye witness to the events depicted. This is what Evola refers to as clairvoyance. Yet, the important thing is not to deal with the external facts surrounding the symbol, but to go deeper into the consciousness evoked by the symbol. This involves penetrating in succession the vital, mental, and spiritual layers. These layers are related to the Degrees of Knowledge\footnote{\url{https://gornahoor.net/?p=411}}.

So, the spiritual layer, or \emph{nous}, is what Tomberg calls ``intellectual experience". Evola claims the \emph{nous} is what he means by the solar or Olympian ``spiritual race". Even beyond that is the Primordial State, the Tao, the Unmoved Mover.

The following meditation is based on the use of the Genesis myth of Adam to penetrate to the awareness of the primordial state, and was intended for students studying Meditations on the Tarot. Again, the important thing is not the precise historical factuality, but rather the awareness of deeper and deeper levels of consciousness. This can be done, too, with other symbols. For example, take any of the many myths and symbols mentioned by Evola and Guenon, and do this sort of active meditation. That is how to reach an understanding of the Golden Age, or the Hyperborean origins, for example.

\paragraph{Guided Meditation}
Invite everyone to take the appropriate role — Adam or Eve — the end result will be the same in any case.

\begin{description}
\item[Step 1: ]

Read the relevant passages or look at the symbols 

\item[Step 2: ]

relax, in preparation, concentration without effort (this is not a history test) 

\item[Step 3: ]

Quietly observe the flow of thoughts, without getting drawn in 

\item[Step 4: ]

Imagine what the Garden is like sensually and physically. The beauty, the abundance, the songs of the birds, the animals, the feel of the grass underfoot, the warmth of the sun, the refreshing coolness of the night. No mosquitos, no sunburn, and so on. Visualize the Tree of Life. Then the Tree of the Knowledge of Good and Evil — is it attracting, or repellent, or more like a void? Allow your imagination free reign to visualize the appearance of the Garden. 

\item[Step 5: ]

Move a little deeper. What is life in the Garden like? The experience of pleasure without pain, no hunger, no thirst. How do you experience your mate? Are you attracted by Adam's or Eve's beauty? Do you experience sexual desire? How do you react to nakedness? Again, give free reign to the imagination, always attentive to higher inspiration. 

\item[Step 6: ]

What is the mental life of Adam and Eve like – their feelings, their thoughts? What is it like to have no fear of death, illness, poverty? No anxiety, since there is perfect manifestation. But is there still something lacking? What is the attraction to the Tree of Knowledge? Why is there even the possibility of temptation? What were they thinking? Does that attraction to worldliness still exist? How to overcome it? 

\item[Step 7: ]

What is the spiritual life like of Adam and Eve? Move beyond words and thoughts and try to experience this directly. They are in constant awareness of the presence of God, not as an old man walking around, but in their hearts. What is that like? There is no crisis of faith, no doubts, no feelings of forsakenness and abandonment because that presence is always felt. 

\item[Step 8: ]

Gradually stop the active participation in the imagination. Detach and observe, and in the silence be attentive to any inspiration from the superconsciousness. 

\item[Step 9: ]

As you let go, and in silence contemplate the flow of thoughts, images, feeling, become aware of the primordial state that precedes all thought, images, \& feelings. 

\end{description}


\flrightit{Posted on 2010-08-13 by Cologero }

\begin{center}* * *\end{center}

\begin{footnotesize}\begin{sffamily}



\texttt{Angel Andre on 2010-10-23 at 14:28 said: }

You speak of Hermetic meditation, most interesting since it is true i feel, most definitely, that without practice one cannot truly know or experience these things. Im sure you have heard of this before, but you should post some of this up and have a discussion on it, id be interested to see what comes of it. Here is what i think is to be one of the most fascinating books in regards to not just theories, but actual doing, actual practice. 

About the man :

\url{http://en.wikipedia.org/wiki/Franz\_Bardon}

some accounts of his life, written by him but also finished by his secretary at the time.

\url{http://www.scribd.com/doc/301829/Frabato-The-Magician-}

And, here is the work, the theory and practice…Initiation Into Hermetics:

\url{http://www.scribd.com/doc/301820/Initiation-into-Hermetics-by-Franz-Bardon}


\end{sffamily}\end{footnotesize}
